\documentclass[11pt,a4paper]{article}

\usepackage[T1]{fontenc}
\usepackage{lmodern}
\usepackage[margin=21mm,headheight=15pt]{geometry}
\usepackage{amsmath,amssymb,mathtools}
\usepackage{microtype}
\usepackage{xcolor}
\usepackage{enumitem}
\usepackage{fancyhdr}
\usepackage{titlesec}
\usepackage{needspace}

\definecolor{examblue}{HTML}{17365D}
\definecolor{answergreen}{HTML}{245C3A}
\definecolor{rulegray}{HTML}{D6DCE4}

\pagestyle{fancy}
\fancyhf{}
\lhead{\textsc{Machine Learning 1}}
\rhead{\textsc{Mock Exam A -- Answer Key}}
\cfoot{\thepage}
\renewcommand{\headrulewidth}{0.4pt}

\titleformat{\section}
  {\Large\bfseries\color{examblue}}
  {}{0pt}{}
  [\vspace{0.15em}\color{rulegray}\titlerule]
\titleformat{\subsection}
  {\normalsize\bfseries\color{answergreen}}
  {}{0pt}{}

\setlength{\parindent}{0pt}
\setlength{\parskip}{0.4em}
\setlist[enumerate]{leftmargin=1.8em,itemsep=0.8em,topsep=0.35em}
\allowdisplaybreaks

\newcommand{\exercise}[2]{%
  \Needspace{8\baselineskip}%
  \par\addvspace{1.1em}%
  {\Large\bfseries\color{examblue}Exercise #1: #2\par}%
  \nobreak\vspace{0.2em}%
  {\color{rulegray}\hrule height 0.5pt}%
  \vspace{0.55em}%
}
\newcommand{\answer}{\textbf{Answer.}\ }

\begin{document}

\begin{center}
  {\Huge\bfseries\color{examblue} Machine Learning 1}\par
  \vspace{0.35em}
  {\LARGE\bfseries Mock Exam A -- Answer Key}\par
  \vspace{0.7em}
  {\large Worked solutions for all 100 points}
\end{center}

\vspace{0.8em}
\fcolorbox{rulegray}{white}{%
  \begin{minipage}{0.94\textwidth}
    Equivalent algebraic forms are valid. The derivations below use numerator-layout derivatives where a scalar derivative is a row vector.
  \end{minipage}%
}

\exercise{1}{Multivariate derivatives and regularization}

\subsection*{1. Temperature-scaled log-softmax}

\answer Since
\[
  \ell_i(x)=\frac{x_i}{\tau}
  -\log\left(\sum_{k=1}^{d}e^{x_k/\tau}\right),
\]
we obtain
\begin{align*}
  \frac{\partial\ell_i}{\partial x_j}
  &=\frac{\delta_{ij}}{\tau}
    -\frac{e^{x_j/\tau}/\tau}{\sum_{k=1}^{d}e^{x_k/\tau}}\\
  &=\frac{1}{\tau}(\delta_{ij}-s_j).
\end{align*}

\subsection*{2. Full Jacobian}

\answer Row $i$, column $j$ must contain $(\delta_{ij}-s_j)/\tau$. Therefore
\[
  \boxed{J_{\ell}(x)=\frac{1}{\tau}
  \left(I_d-\mathbf{1}s(x)^\top\right)}
  \in\mathbb{R}^{d\times d}.
\]
The matrix $\mathbf{1}s^\top$ has entry $s_j$ in every row.

\subsection*{3. Invariant direction}

\answer Since $s^\top\mathbf{1}=1$,
\[
  J_{\ell}(x)\mathbf{1}
  =\frac{1}{\tau}\left(\mathbf{1}
  -\mathbf{1}s^\top\mathbf{1}\right)=0.
\]
For any scalar $c$,
\[
  s_i(x+c\mathbf{1})
  =\frac{e^{c/\tau}e^{x_i/\tau}}
  {e^{c/\tau}\sum_k e^{x_k/\tau}}
  =s_i(x).
\]
Thus $\ell(x+c\mathbf{1})=\ell(x)$. The directional derivative along $\mathbf{1}$ is consequently zero.

\Needspace{16\baselineskip}
\subsection*{4. Generalized ridge regression}

\answer The numerator-layout derivative is
\[
  \frac{\partial L}{\partial\theta}
  =\frac{2}{n}(X\theta-y)^\top X+2\gamma\theta^\top R.
\]
After transposing the stationarity equation,
\[
  \frac{1}{n}X^\top(X\theta-y)+\gamma R\theta=0,
\]
so
\[
  (X^\top X+n\gamma R)\theta=X^\top y,
  \qquad
  \boxed{\widehat{\theta}
  =(X^\top X+n\gamma R)^{-1}X^\top y}.
\]
For any nonzero $v\in\mathbb{R}^p$,
\[
  v^\top(X^\top X+n\gamma R)v
  =\lVert Xv\rVert_2^2+n\gamma v^\top Rv>0.
\]
Thus the matrix is positive definite even if $X$ is rank deficient. Equivalently, the Hessian
\[
  \frac{2}{n}X^\top X+2\gamma R
\]
is positive definite. Hence $L$ is strictly convex and the stationary point is the unique global minimizer.

\exercise{2}{Waiting times and Bayesian estimation}

\subsection*{1. Expected waiting time}

\answer Put $r=1-q$. Differentiating the geometric series gives
\[
  \sum_{k=1}^{\infty}kr^{k-1}=\frac{1}{(1-r)^2},
  \qquad
  \sum_{k=0}^{\infty}kr^k=\frac{r}{(1-r)^2}.
\]
Therefore
\[
  \mathbb{E}[K\mid q]
  =q\sum_{k=0}^{\infty}k(1-q)^k
  =q\frac{1-q}{q^2}
  =\boxed{\frac{1-q}{q}}.
\]

\subsection*{2. Requested probability}

\answer
\[
  P(K=2\mid q=0.25)=0.25(0.75)^2
  =\boxed{\frac{9}{64}=0.140625}.
\]

\subsection*{3. Likelihood and log-likelihood}

\answer By independence,
\[
  L(q)=\prod_{n=1}^{N}q(1-q)^{K_n}
  =q^N(1-q)^S.
\]
Thus
\[
  \boxed{\log L(q)=N\log q+S\log(1-q)}
\]
up to additive constants independent of $q$.

\subsection*{4. Maximum-likelihood estimator}

\answer
\[
  \frac{d}{dq}\log L(q)=\frac{N}{q}-\frac{S}{1-q}=0
  \quad\Longrightarrow\quad
  \boxed{\widehat q_{\mathrm{ML}}=\frac{N}{N+S}}.
\]
Because $S>0$, this value is in $(0,1)$. Also,
\[
  \frac{d^2}{dq^2}\log L(q)
  =-\frac{N}{q^2}-\frac{S}{(1-q)^2}<0,
\]
so the log-likelihood is strictly concave and this is its unique maximum.

\subsection*{5. Posterior distribution}

\answer Multiplication of the likelihood and prior gives
\begin{align*}
  p(q\mid K_1,\ldots,K_N)
  &\propto q^N(1-q)^S q^{a-1}(1-q)^{b-1}\\
  &=q^{a+N-1}(1-q)^{b+S-1}.
\end{align*}
Therefore
\[
  \boxed{q\mid K_1,\ldots,K_N
  \sim\operatorname{Beta}(a+N,b+S)}.
\]

\subsection*{6. MAP estimator}

\answer The derivative of the posterior log-density is
\[
  \frac{a+N-1}{q}-\frac{b+S-1}{1-q}.
\]
Setting it to zero gives
\[
  \boxed{\widehat q_{\mathrm{MAP}}
  =\frac{a+N-1}{a+b+N+S-2}}.
\]
Both posterior shape parameters exceed one, so the mode is unique and interior.

\subsection*{7. Data statistic}

\answer The statistic is
\[
  \boxed{S=\sum_{n=1}^{N}K_n}.
\]
For fixed $N$, both the likelihood and posterior depend on the observed values only through $S$. Thus $S$ is sufficient for $q$ by the factorization criterion.

\exercise{3}{Estimating an unknown range}

\subsection*{1. Likelihood}

\answer Each observation has density
\[
  p(x_n\mid M)=\frac{1}{2M}\mathbf{1}\{|x_n|\leq M\}.
\]
Consequently,
\[
  \boxed{p(x_1,\ldots,x_N\mid M)
  =(2M)^{-N}\mathbf{1}\{M\geq A\}},
  \qquad M>0.
\]

\subsection*{2. Maximum-likelihood estimator}

\answer The likelihood is zero for $M<A$ and equals $(2M)^{-N}$ for $M\geq A$. It is strictly decreasing over the feasible range, so
\[
  \boxed{\widehat M_{\mathrm{ML}}=A=\max_n|X_n|}.
\]
For the supplied observations,
\[
  A=\max(3.2,1.1,2.4,2.8,0.5)=\boxed{3.2}.
\]

\subsection*{3. Unbiased correction}

\answer For $\widetilde M=cA$, unbiasedness requires
\[
  \mathbb{E}[\widetilde M]
  =c\frac{NM}{N+1}=M.
\]
Thus
\[
  \boxed{\widetilde M=\frac{N+1}{N}A}.
\]

\subsection*{4. Variance}

\answer First,
\begin{align*}
  \operatorname{Var}(A)
  &=\frac{NM^2}{N+2}-\frac{N^2M^2}{(N+1)^2}\\
  &=\frac{NM^2}{(N+2)(N+1)^2}.
\end{align*}
Therefore
\[
  \boxed{\operatorname{Var}(\widetilde M)
  =\left(\frac{N+1}{N}\right)^2\operatorname{Var}(A)
  =\frac{M^2}{N(N+2)}}.
\]

\exercise{4}{Multi-output regression with correlated noise}

\subsection*{1. Dimensions}

\answer
\[
  \boxed{\Phi W\in\mathbb{R}^{N\times k}},\qquad
  \boxed{T-\Phi W\in\mathbb{R}^{N\times k}},\qquad
  \boxed{(T-\Phi W)^\top(T-\Phi W)\in\mathbb{R}^{k\times k}}.
\]

\subsection*{2. Log-likelihood}

\answer Let $E(W)=T-\Phi W$. Since
\[
  |\sigma^2C|=(\sigma^2)^k|C|,
\]
independence gives
\begin{align*}
  \log p(T\mid W,\sigma^2)
  ={}&-\frac{Nk}{2}\log(2\pi)
  -\frac{Nk}{2}\log\sigma^2
  -\frac{N}{2}\log|C|\\
  &-\frac{1}{2\sigma^2}
  \operatorname{tr}\left(C^{-1}E(W)^\top E(W)\right).
\end{align*}

\Needspace{16\baselineskip}
\subsection*{3. Maximum-likelihood estimator of $W$}

\answer For fixed $\sigma^2$, minimize
\[
  Q(W)=\operatorname{tr}(C^{-1}E^\top E).
\]
Using $dE=-\Phi\,dW$ and symmetry of $C^{-1}$,
\[
  dQ=-2\operatorname{tr}(C^{-1}E^\top\Phi\,dW),
\]
so stationarity gives
\[
  \Phi^\top(T-\Phi W)C^{-1}=0.
\]
Since $C$ is invertible,
\[
  \Phi^\top\Phi W=\Phi^\top T,
  \qquad
  \boxed{W_{\mathrm{ML}}
  =(\Phi^\top\Phi)^{-1}\Phi^\top T}.
\]
The positive scalar $1/\sigma^2$ cancels, and multiplication by the invertible $C^{-1}$ does not change the normal equations. Hence the estimate depends on neither $C$ nor $\sigma^2$.

\subsection*{4. Maximum-likelihood estimator of $\sigma^2$}

\answer Let $E=T-\Phi W_{\mathrm{ML}}$ and
\[
  Q=\operatorname{tr}(C^{-1}E^\top E)>0.
\]
For $z=\sigma^2$, the relevant log-likelihood is
\[
  -\frac{Nk}{2}\log z-\frac{Q}{2z}.
\]
Its derivative is $-Nk/(2z)+Q/(2z^2)$. Therefore
\[
  \boxed{\widehat\sigma_{\mathrm{ML}}^2
  =\frac{\operatorname{tr}(C^{-1}E^\top E)}{Nk}}.
\]
At this value the second derivative is $-Nk/(2z^2)<0$. The nonzero-residual assumption ensures a positive estimate and an attained maximum.

\subsection*{5. Residual orthogonality}

\answer The normal equations imply
\[
  \boxed{\Phi^\top E
  =\Phi^\top(T-\Phi W_{\mathrm{ML}})=0}.
\]
Thus every residual column is orthogonal to every column of $\Phi$. For each output, the fitted vector is the Euclidean orthogonal projection of the target onto the column space of $\Phi$.

\exercise{5}{Information from one Bayesian observation}

Write
\[
  D=\beta^{-1}+\phi^\top\Sigma_N\phi>0.
\]

\subsection*{1. Covariance update}

\answer Apply Sherman--Morrison to
\[
  \Sigma_{N+1}=(\Sigma_N^{-1}+\beta\phi\phi^\top)^{-1}.
\]
This gives
\begin{align*}
  \Sigma_{N+1}
  &=\Sigma_N-
  \frac{\beta\Sigma_N\phi\phi^\top\Sigma_N}
  {1+\beta\phi^\top\Sigma_N\phi}\\
  &=\boxed{\Sigma_N-
  \frac{\Sigma_N\phi\phi^\top\Sigma_N}
  {\beta^{-1}+\phi^\top\Sigma_N\phi}}.
\end{align*}

\subsection*{2. Innovation form of the mean}

\answer Start from
\[
  \mu_{N+1}=\Sigma_{N+1}
  (\Sigma_N^{-1}\mu_N+\beta\phi t).
\]
The covariance update implies
\[
  \Sigma_{N+1}\Sigma_N^{-1}
  =I-\frac{\Sigma_N\phi\phi^\top}{D},
  \qquad
  \Sigma_{N+1}\beta\phi=\frac{\Sigma_N\phi}{D}.
\]
Hence
\[
  \boxed{\mu_{N+1}=\mu_N+K(t-\phi^\top\mu_N)},
  \qquad
  \boxed{K=\frac{\Sigma_N\phi}{D}}.
\]
The vector $\Sigma_N\phi$ selects the parameter direction correlated with the measured quantity $\phi^\top w$. The scalar innovation $t-\phi^\top\mu_N$ gives the signed prediction error and controls the update magnitude.

\subsection*{3. Posterior predictive distribution}

\answer Write $t=\phi^\top w+\varepsilon$, where $w\mid\mathcal D_N\sim\mathcal N(\mu_N,\Sigma_N)$ and $\varepsilon\sim\mathcal N(0,\beta^{-1})$ are independent. Therefore
\[
  \boxed{p(t\mid\phi,\mathcal D_N)
  =\mathcal N\left(t\,\middle|\,
  \phi^\top\mu_N,
  \beta^{-1}+\phi^\top\Sigma_N\phi\right)}.
\]
The variance consists of observation-noise uncertainty $\beta^{-1}$ and posterior parameter uncertainty $\phi^\top\Sigma_N\phi$.

\subsection*{4. Directional uncertainty reduction}

\answer From the covariance update,
\begin{align*}
  a^\top(\Sigma_N-\Sigma_{N+1})a
  &=\frac{a^\top\Sigma_N\phi\phi^\top\Sigma_Na}{D}\\
  &=\boxed{\frac{(a^\top\Sigma_N\phi)^2}
  {\beta^{-1}+\phi^\top\Sigma_N\phi}}\geq 0.
\end{align*}
It is zero exactly when $a^\top\Sigma_N\phi=0$. This quantity is the posterior covariance between $a^\top w$ and $\phi^\top w$. Thus the observation reduces uncertainty only in parameter directions correlated with the measured quantity; the rank-one reduction acts along $\Sigma_N\phi$.

\vfill
\begin{center}
  \textbf{End of Mock Exam A Answer Key}
\end{center}

\end{document}
